\section{Conjugacy, Centralisers and Normalisers}

A group \(G\) can act on itself by conjugation as
\[
    \function{\ast}{G \times G}{G}{\pqty{g, h}}{g h g\inverse}.
\]

\begin{lemma}[Conjugation by Element gives Homomorphism from the Group to Itself]
    \label{lem:conjugation-element-homomorphism-self}%
    If we fix \(g \in G\), the permutation
    \[
        \function{\pi}{G}{G}{h}{ghg\inverse}
    \]
    gives a homomorphism from \(G\) to itself.
\end{lemma}

\begin{definition}[Automorphism, Automorphism Group]
    \label{def:automorphism}%
    Let \(G\) be a group. A permutation \(G \to G\) is an \emph{automorphism} if it is also a homomorphism.

    The set of all automorphisms of \(G\) form a group, namely the \emph{automorphism group}, denoted as
    \[
        \Aut\pqty{G} = \set{\functiondomain{f}{G}{G}}{f \text{ is an automorphism}} \subgroup \Sym\pqty{G}.
    \]
\end{definition}

\begin{example*}
    Let \(G = \RR^2\) be under addition. Then \(\Aut\pqty{G} \supgroup \GL_2\pqty{\RR}\).
\end{example*}

\begin{notation}
    The automorphisms given by conjugation action are called \emph{inner automorphisms}, denoted as \(\Inn\pqty{G}\). This is a normal subgroup of \(\Aut\pqty{G}\), and the quotient is the \emph{outer automorphism group}, denoted as \(\Out\pqty{G}\).
\end{notation}

\begin{definition}[Conjugation Class, Centraliser, Centre]
    \label{def:ccl-centraliser-centre}%
    Fix \(g \in G\). The \emph{conjugacy class} of \(g\) is the orbit under conjugation action, given by
    \[
        \ccl_G\pqty{g} = \set{hgh\inverse}{h \in G}.
    \]

    The \emph{centraliser} of \(g \in G\) is the stabiliser under conjugation action, given by
    \[
        C_G\pqty{g} = \set{h \in G}{h g h\inverse = g}.
    \]

    The \emph{centre} of \(G\) is the kernel of conjugation action, given by
    \[
        Z\pqty{G} = \set{z \in G}{\forall g \in G: zg = gz}.
    \]
\end{definition}

\begin{proposition}[Size of Conjugacy Class]
    \label{prop:size-conjugacy-class}%
    Let \(G\) be a finite group. Then for all \(x \in G\), we have
    \[
        \abs{\ccl\pqty{x}} = \abs{G : C_G\pqty{x}} = \frac{\abs{G}}{\abs{C_G\pqty{x}}}.
    \]
\end{proposition}

\begin{proof}
    Apply \autoref{thm:orbit-stabiliser} (Orbit--Stabiliser Theorem).
\end{proof}

\begin{definition}[Normaliser]
    \label{def:normaliser}%
    Let \(H \subgroup G\). The \emph{normaliser} of \(H \in G\) is
    \[
        N_G \pqty{H} = \set{g \in G}{gHg\inverse = H}.
    \]
\end{definition}

Observe that \(H \normalsub N_G\pqty{H}\), and in fact by definition, \(N_G\pqty{H}\) is the \emph{largest} subgroup of \(G\), inside which \(H\) is normal.